\section{Introduction}
\label{sec:intro}

Flaky tests pass and fail unpredictably on the same commit of production code. They erode developer trust in continuous integration (CI), stall release pipelines, and mask genuine regressions~\cite{luo2014empirical,parry2021survey}. Among non-concurrency root causes, two categories dominate automated flaky test repair across Java and Python suites. \emph{Implementation-Dependent} (ID) tests rely on unspecified iteration orders in hash collections, reflection APIs, or JSON serializers~\cite{gyori2016nondex}. \emph{Order-Dependent} (OD) tests depend on shared mutable state polluted or initialized by earlier tests~\cite{lam2019idflakies,shi2019ifixflakies,wang2022ipflakies}.

Automated repair for these tests has advanced rapidly from suite search and syntactic templates, such as \textsc{iFixFlakies}~\cite{shi2019ifixflakies}, \textsc{DexFix}~\cite{zhang2021dexfix}, and \textsc{iPFlakies}~\cite{wang2022ipflakies}, to iterative large language model (LLM) agents including \textsc{FlakyDoctor}~\cite{chen2024flakydoctor} and industrial repair pipelines~\cite{rastenis2026autocover,leunig2026perl}. Across these systems, however, patch acceptance remains governed by a \emph{single-sided stabilization validator}. Once a candidate edit compiles, passes \textsc{NonDex}~\cite{gyori2016nondex} or order-schedule reruns, and keeps its assertion statements intact, the pipeline accepts the patch and stops.

\paragraph{Stabilization via Oracle Erosion.}
Optimizing solely for execution stability creates a subtle quality hazard that we call \emph{Stabilization via Oracle Erosion}. In automated program repair on production code~\cite{qi2015kali,smith2015cure}, an external test suite constrains candidate edits and exposes broken functionality. Flaky test repair lacks that safeguard because it modifies the test oracle and fixture themselves. A concrete example from Apache Commons Lang illustrates the mechanism. In \texttt{ToStringBuilderTest}, an LLM repair replaced a strict reflection-output check with Java's single-argument \texttt{split(",")} to remove iteration-order sensitivity. That patch passed every stability rerun, but single-argument splitting silently drops trailing empty tokens and left eight production faults in \texttt{ToStringStyle} undetected. Once merged, an eroded test turns permanently green and raises no CI alarm.

Warning signs of this hazard have appeared across separate lines of software quality research. \textsc{TestBoost}~\cite{li2024testboost} found that nearly a third of fixture-transformed test classes lost mutation-killing ability despite passing order checks. In this journal, Leunig and Petersen~\cite{leunig2026perl} reported that LLM-driven test repair at IBM frequently diluted assertion specificity behind green pass rates, and Bludau and Pretschner~\cite{bludau2026szz} showed how noisy evaluation proxies distort reported effectiveness. Four empirical questions remain unanswered. First, once we control for flaky mutants~\cite{parry2026flimsiness} and verify actual patch provenance, how often does silent oracle erosion occur across historical tools and modern LLMs, and what false-rejection boundaries emerge on developer pull requests? Second, which edit mechanisms drive erosion, and why do stronger static structural guards fail to block them? Third, how does a fail-closed, test-scoped mutation gate balance operational safety against abstention and held-out yield, and how much recovery do localized witnesses add over bare numeric counts? Fourth, what uncached probing times, abstention causes, and mutant-subset trade-offs arise in practice?

\paragraph{Cross-Language Study and MuGuard-Flake.}
To measure erosion without grading a gate on the same mutants it optimizes against, we build a dual-language benchmark of $\BenchEligible$ flaky tests across $\BenchProjects$ Java and Python repositories. Each baseline fault obligation is verified across repeated clean runs in the test's dynamically executed focal scope, alongside separate obligations for shared-state partner tests whenever order-dependent edits touch fixtures (Section~\ref{sec:construct}). Before inspecting any candidate repair, we freeze $\BenchMExtTotal$ disjoint held-out faults and independent stability schedules to evaluate generalization cleanly. We also collect and manually audit $\PRTotalInstances$ PR-linked Java test instances to measure false rejections on legitimate human fixes.

Building on this benchmark, we design \textsc{MuGuard-Flake}, a fail-closed local mutation gate inspired by error-controlling LLM testing gates~\cite{salva2026nltestrunner} and diagnostic feedback studies~\cite{zhang2026typerepair} (Section~\ref{sec:methodology}). Any candidate failing compilation or stability reruns short-circuits immediately before mutation probing. For stabilized patches, \textsc{MuGuard-Flake} maps every edited test and affected partner scope, probes baseline focal obligations in seconds, and feeds up to four surviving-mutant witnesses back to the model on erosion retries. If a five-round budget expires without an obligation-preserving fix, the gate abstains and hands those witnesses to maintainers instead of merging a weakened test.

\paragraph{Contributions.}
This paper makes four empirical and engineering contributions to software quality assurance:
\begin{enumerate}
    \item \emph{Cross-Language Erosion Audit and Pull-Request Boundaries.} Across $\BenchEligible$ flaky tests, historical repair tools, and five LLMs from four vendors, roughly one in five to one in four stabilization-passing patches silently drops baseline obligations. Checking both target and partner scope across manually reviewed developer pull requests establishes a \PRLegitFalseRejMuGuardPct{} false-rejection rate on legitimate repairs alongside subtle weakening patches missed by standard mutators.
    \item \emph{Edit-Mechanism Taxonomy and Static-Guard Blind Spots.} Auditing \TaxUniquePatches{} unique eroded patches alongside a \TaxManualN-patch manual review explains why over four-fifths of eroded patches bypass an enhanced static structural guard. Precondition shifts, delimiter normalizations, partner-scope edits, and focal call bypasses evade syntactic rules, and static heuristics even reduce sound-fix yield on an out-of-sample model.
    \item \emph{Fail-Closed Mutation Gating and Witness Isolation.} \textsc{MuGuard-Flake} eliminates erosion during repair across all five models, raises operational sound-fix yield by 12 to 18 percentage points over static guards on multi-round backbones, and cuts unseen held-out fault loss among accepted patches by 45\% to 65\%. Conservative abstention leaves overall end-to-end held-out yield statistically inconclusive, whereas a controlled feedback comparison shows that concrete surviving-mutant witnesses recover far more eroded repairs than bare numeric counts.
    \item \emph{Uncached Probing Budgets and Open Replication Package.} Even without caching, focal mutation probing requires under ten seconds median in Java and under four seconds in Python, whereas model reasoning dominates end-to-end latency. We diagnose why abstentions occur, quantify the safety--cost trade-off of sampling smaller mutant subsets, and release all manifests, regression tests, execution traces, and verification scripts.
\end{enumerate}

The remainder of this paper reviews related work (Section~\ref{sec:related}), formalizes oracle erosion and external validation (Section~\ref{sec:construct}), details \textsc{MuGuard-Flake} (Section~\ref{sec:methodology}), presents our experimental design and results (Sections~\ref{sec:design}--\ref{sec:results}), and discusses deployment boundaries, examines validity threats, and concludes (Sections~\ref{sec:discussion}--\ref{sec:conclusion}).
